\section{Group Rings}
\setcounter{subsection}{0}

\subsection{The Units Conjecture}
Let $K$ be a field, $G$ a group.
Let $K[G]$ denote the group algebra with coefficients in $K$.

The \textbf{Units Conjecture} proposed the following:
\textit{If $K$ is a field and $G$ is a torsion-free group,
then all units in $K[G]$ are trivial.}
It was proven false in 2021 by Giles Gardam.

\begin{exr}
	Let $V=\bZ_2\times \bZ_2$ be the Klein $4$-group.
	Show that the group algebra $\bZ_2[V]$ has a nontrivial unit.
\end{exr} \

\begin{soln}
	The element $1\cdot(1,1)$ is a unit, with itself as its own inverse.
	To see this, we can write it as:
	\begin{equation*}
		1\cdot(1,1) \ = \ 0\cdot(0,0)+0\cdot(1,0)+0\cdot(0,1)+1\cdot(1,1)
	\end{equation*}
	Squaring this expression gives our result by using the intuitive
	multiplication.
\end{soln}

\newpage
\begin{exr}
	Let $G$ be a torsion-free group.
	Prove that if $K[G]$ has only trivial units,
	then it has no nonzero nilpotents.
\end{exr} \

\begin{pf}
	Suppose $a^{n}=0$ for some $a\in K[G]$.
	Then, notice:
	\begin{equation*}
		(1-a)(a+a^{2}+\dots +a^{n-1})
		\ = \ 1-a^{n} \ = \ 1
	\end{equation*}
	Thus, $1-a$ is a unit.
	Since all units of $K[G]$ are trivial, then $1-a=kg$ for some $k\in K,g\in G$.
	But if $a^{n}=(1-kg)^{n}=0$, then by considering its binomial expansion,
	we must have that $k=0$ or $g=e$.
	If $k=0$, then $1-a$ is not a unit; if $g=e$, then $a=0$.
\end{pf} \

\begin{exr}
	Suppose that $K[G]$ has only trivial units.
	Suppose that $|K|,|G|\geq 3$. Prove that $K[G]$ is semiprimitive.
\end{exr} \

\begin{pf}
	If $a\in J(K[G])$, then $1-a=kg$ is a unit.
	If $g=e$, then $a=1-k\in K$, but $a\in J(K[G])$ implies $1\in J(K[G])$,
	a contradiction.
	If $g\neq e$, then since $|K|\geq3$, pick $c\neq0,1$.
	Then $1-ca=(1-c)+ckg$ is a unit, so necessarily $a=0$.
\end{pf}


\newpage
\subsection{Dietzmann's Lemma}
Stating, but not proving.

\begin{lm}
	Let $G$ be a group and let $Z(G)$ be its center.
	Prove that if $[G:Z(G)]<\infty$,
	then the commutator subgroup $[G,G]$ is finite.
\end{lm}


\newpage
\subsection{Connell's Theorem - Part I}
Let $G$ be a group.
Let $\Delta(G)$ be the set of all elements $x\in G$ whose conjugacy class is finite ---
equivalently, the centralizer $C_G(x)$ has finite index in $G$ ---
and let $\Delta^+(G)$ be the set of all elements of $\Delta(G)$ of finite order.
In this problem we will show that $\Delta^+(G)$ is a subgroup of $G$.

\begin{exr}
	Let $g_1,\ldots,g_n$ be elements of $\Delta^+(G)$.
	Let $N$ be the subgroup generated by all the $g_i$'s and their conjugates.
	Show that $N$ is a finite normal subgroup of $G$.

	Hint: Use Dietzmann's lemma.
\end{exr} \

\begin{exr}
	Prove that $\Delta^+(G)$ is a subgroup of $G$.
\end{exr} \

\begin{exr}
	Prove that $G$ has no finite normal subgroups if and only if $\Delta^+(G)=\{1\}$.
\end{exr}


\newpage
\subsection{Connell's Theorem - Part II}
Our goal is to show that if $\Delta^+(G)=\{1\}$, then $K[G]$ is a prime ring.
Here we show that $K[\Delta(G)]$ is a prime ring.

\begin{exr}
	Let $T$ be a torsion-free group and assume that
	every element of $T$ has finite conjugacy class.
	Prove that $K[T]$ is a commutative domain.
\end{exr} \

\begin{exr}
	Suppose that $G$ has no finite normal subgroups.
	Prove that $K[\Delta(G)]$ is a commutative domain.

	Hint: Apply part (a) to $T=\Delta(G)$.
\end{exr}


\newpage
\subsection{Neumann's Lemma}
Stating, but not proving.

\begin{lm}
	Let $G$ be a group with subgroups $H_1,\ldots,H_m$.
	Suppose that $G$ can be expressed as a finite union of
	cosets of the $H_i$'s; that is:
	\begin{equation*}
		G = \bigcup_{i=1}^m\bigcup_{j=1}^{n_i} H_i a_{ij}
	\end{equation*}
	Prove that at least one of the $H_i$'s has finite index in $G$.
\end{lm}


\newpage
\subsection{Connell's Theorem - Part III}
Now we know that $K[\Delta(G)]$ is a (commutative) prime ring.
Now we show that this implies that $K[G]$ is also prime.

Let $H$ be a subgroup of $G$.
Then $K[H]$ is a subring of $K[G]$.
Define the \textbf{restriction map} $\pi_H:K[G]\rightarrow K[H]$ by $\pi(\alpha):=\alpha\vert_H$.
For a right ideal $I$ in $K[H]$, let $\rho(I)$ be the right ideal of $K[G]$ generated by $I$.

\begin{exr}
	Prove that $I\subseteq \rho(\pi(I))$.
	Thus show that $I=0$ if and only if $\pi(I)=0$.
\end{exr} \

\begin{exr}
	Let $I,J$ be two ideals in $K[G]$.
	Suppose that $IJ=0$. Prove that $\pi(I)\pi(J)=0$. 

	Hint: This is tricky. Let $x\in I$ and $y\in J$. We must show that $\pi(x)\pi(y)=0$.
	First convince yourself it is enough to show that $\pi(x)y=0$.
	Now, $x$ can be expressed as a linear combination of elements of $G$;
	split this up as $x=x_0+x_1$ is the part in $K[\Delta(G)]$, and $x_1$ is the rest (thus $\pi(x)=x_0$).
	Say $x_0=\sum_{i=1}^ma_iu_i$ where $u_i\in \Delta(G)$, and
	$x_1=\sum_{j=1}^nb_jv_j$, where $v_j\not\in \Delta(G)$.
	Show that the intersection of centralizers $C=\bigcap_{i=1}^m C_G(u_i)$
	has finite index in $G$, and that (if $x_0y\neq 0$) then it is contained in
	a union of the cosets of the $C_G(v_i)$'s.
	Now use Neumann's Lemma.
\end{exr} \

\begin{exr}
	Let $G$ be any group.
	Show that if $K[\Delta(G)]$ is prime, then so is $K[G]$.
\end{exr} \

\begin{exr}
	Deduce Connell's Theorem: that is,
	if $K$ is a field and $G$ is a group, then $K[G]$ is prime
	if and only if $\Delta^+(G)=\{1\}$.
\end{exr}


\newpage
\subsection{Faithful Trace Theorem - Part I}
A ring $R$ is said to be \textbf{Dedekind-finite} if $ab=1\implies ba=1$ for all $a,b\in R$.
Our goal is to prove that the group ring $\bC[G]$ is always Dedekind-finite, for any group $G$.
To do this, we show that $\bC[G]$ possesses a ``faithful trace'';
from this it is easy to see that every algebra possessing a faithful trace is necessarily Dedekind-finite.

Let $R$ be a $\bC$-algebra.
A \textbf{trace} on $R$ is a linear functional $\tau:R\rightarrow \bC$ such that
$\tau(1)=1$ and $\tau(ab)=\tau(ba)$ for all $a,b\in R$.
A trace $\tau$ is said to be \textbf{faithful} if $\tau(e)=0\implies e=0$ whenever
$e=e^2$ is an idempotent in $R$.

\begin{exr}
	Prove that if $R$ has a faithful trace, then $R$ is Dedekind-finite.
\end{exr} \

\begin{pf}
	Suppose $ab=1$ and $\tau$ is a faithful trace on $R$.
	Then $1-ba$ is an idempotent:
	\begin{equation*}
		(1-ba)^{2} \ = \ 1-2ba+(ba)^{2} \ = \ 1-2ba+(ba)(ba) \ = \ 1-2ba+ba \ = \ 1-ba
	\end{equation*}
	It also has zero trace:
	\begin{equation*}
		\tau(1-ba) \ = \ \tau(1)-\tau(ba) \ = \ \tau(1)-\tau(ab) \ = \ 1-1 \ = \ 0
	\end{equation*}
	Thus, $1-ba=0$, so $ba=1$ as needed.
\end{pf} \

\begin{exr}
	Prove that the normalized trace on $M_n(\bC)$ is faithful.
\end{exr} \

\begin{pf}
	Recall that $\trm{tr}(A)$ is the sum of the eigenvalues of $A$ with multiplicity.
	If $A^{2}=A$, then $A$ is a projection, so $\trm{tr}(A)=\trm{rank}(A)$,
	and the result follows.
\end{pf} \

\begin{exr}
	For a group $G$, define a function $\tau$ on $\bC[G]$ by $\tau(\alpha):=\dfrac{1}{|G|}\alpha(1)$.
	Prove that $\tau$ is a trace on $\bC[G]$.
\end{exr}
ignoring factor cz idt thats generally true.
also doesnt say $G$ is finite whats up w that

\begin{pf}
	Evidently, $\tau(1)=1$. We also have:
	\begin{equation*}
		\tau(\alpha\beta) \ = \ \sum_{g\in G}\alpha(g)\beta(g^{-1})
		\ = \ \sum_{g\in G}\alpha(g^{-1})\beta(g) \ = \ \tau(\beta\alpha)
	\end{equation*}
\end{pf}


\newpage
\subsection{Faithful Trace Theorem - Part II}
Let $G$ be a group.
For an element  $\alpha=\sum_{g\in G}\alpha_gg\in\bC[G]$,
define the \textbf{conjugate} of $\alpha$ by:
\begin{equation*}
	\alpha^* := \sum_{g\in G} \bar{\alpha_g}g^{-1}
\end{equation*}
For $\alpha,\beta\in \bC[G]$, define the \textbf{inner product} of $\alpha$ with $\beta$ by:
\begin{equation*}
	\langle \alpha,\beta\rangle := \tau(\alpha\beta^*)
\end{equation*}

\begin{exr}
	A \textbf{self-adjoint idempotent} in $\bC[G]$ is an element $e$ such that $e=e^2=e^*$.
	Show that if $e$ is a self-adjoint idempotent and $\tau(e)=0$, then $e=0$.

	This shows that faithfulness holds if $e$ is assumed to be self-adjoint.
	Thus our goal is to show that every idempotent can be replaced by a self-adjoint one.
\end{exr} \

\begin{pf}
	We have that:
	\begin{equation*}
		\tau(e) \ = \ \tau(ee) \ = \ \tau(ee^{*}) \ = \ \ang{e,e} \ = \ 0
	\end{equation*}
	Thus, $e=0$ as needed.
\end{pf} \

\begin{exr}
	Show that $\langle\cdot,\cdot\rangle$ is an inner product.
	Thus $\bC[G]$ is an inner product space.
\end{exr} \

\begin{exr}
	Let $I$ be a right ideal of $\bC[G]$.
	Show that the orthogonal complement $I^\perp$ is also a right ideal of $\bC[G]$.
\end{exr} \

\begin{pf}
	Suppose $b\in I^{\perp}$. We want to show that $bc\in I^{\perp}$ for any $c\in\bC[G]$.
	Note:
	\begin{equation*}
		\ang{a,bc} \ = \ \tau(a(bc)^{*}) \ = \ \tau(ac^{*}b^{*})
		\ = \ \ang{ac^{*},b}
	\end{equation*}
	Since $I$ is a right ideal, $ac^{*}\in I$ if $a\in I$.
	Since $b\in I^{\perp}$, then by definition and the above, we get that
	$\ang{a,bc}=0$, so $bc\in I^{\perp}$ as needed.
\end{pf} \

\begin{exr}
	Let $e$ be an idempotent, and let $I$ be the right ideal generated by $e$.
	Let $\gamma$ be the distance from $I$ to $1$, that is:
	\begin{equation*}
		\gamma := \inf\{\|\alpha-1\|:\alpha \in I\}
	\end{equation*}
	Prove that $\gamma>0$.
\end{exr} \

\begin{pf}
	First, note that for any $a=\sum ea_{i}\in I$, we have:
	\begin{equation*}
		a \ = \ \sum ea_{i} \ = \ \sum e^{2}a_{i} \ = \ e\left(\sum ea_{i}\right) \ = \ ea
	\end{equation*}
	since $e$ is an idempotent.
	Now, suppose $\gamma=0$.
	Then there is a sequence $a_{n}\in I$ such that $a_{n}\sto1$.
	By the above, we also have $a_{n}=ea_{n}\sto e$,
	so by uniqueness of limits, $e=1$.
\end{pf} \

\begin{exr}
	Let $\gamma$ be as in the previous part.
	Then, let $(f_n)_{n\geq 1}$ be any sequence in $I$ such that
	$\|f_n-1\|\leq \gamma^2+n^{-4}$.
	\begin{enumerate}[i)]
		\item Prove that $|\|f_n\|^2-\tau(f_n)\|\rightarrow 0$ as $n\rightarrow\infty$.
		\item Prove that $\|f_ne-e\|\rightarrow 0$ as $n\rightarrow \infty$.
	\end{enumerate}
\end{exr} \

\begin{exr}
	Finally, show that $\tau$ is faithful.
	Deduce that $\bC[G]$ is Dedekind-finite.

	(In class I'll show you how to use the $\bC$ case to prove a
	general version that works over any field, in very few lines.)
\end{exr}


\newpage
\subsection{Passmann's Primitivity Theorem}
Let $K$ be a field and let $G$ be a group.
When is the group ring $K[G]$ primitive? Here we give a necessary condition:
if $K[G]$ is primitive and $K$ is an algebraically closed field with $|K|>|G|$, then $\Delta(G)=\{1\}$.

\begin{exr}
	Prove that primitive rings are prime.
	Thus deduce that if $K[G]$ is primitive, then $\Delta^+(G)=\{1\}$.
\end{exr} \

\begin{exr}
	Let $K$ be an algebraically closed field, and let $R$ be a primitive
	$K$-algebra with $|K| > \dim_K(R)$.
	Show that $Z(R)$ is $1$-dimensional, as follows.
	\begin{enumerate}[i)]
		\item Show that $Z(R)$ is a commutative domain.
		\item Let $F$ be the fraction field of $Z(R)$.
			Show that $F/K$ is a field extension of finite degree.
			Conclude $F=K$ by algebraic closure.
		\item Conclude $Z(R)=K$.
			(More precisely, that every element of $Z(R)$ has the form $c\cdot 1$,
			where $c\in \bC$.)
	\end{enumerate}
\end{exr} \

\begin{exr}
	Let $C$ be a finite conjugacy class in $G$, and let $z_C$ be
	the sum of elements of $C$.
	Prove that the center of the group ring $K[G]$ is spanned by the $z_C$'s.
\end{exr} \

\begin{exr}
	Now use 2) and 3) [or b) and c)] to prove Passman's Theorem.
\end{exr}
